

\subsection{Notation}



\begin{longtable}{|c|p{10cm}|}
\hline
\textbf{Symbol/Notation} & \textbf{Description} \\ \hline
\endfirsthead

\hline
\textbf{Symbol/Notation} & \textbf{Description} \\ \hline
\endhead

\hline
\endfoot

\hline
\caption{Table with Notation for Problem Formulation} \\
\endlastfoot

\( T \) & Original tabular data containing both continuous and discrete variables. \\ \hline
\( n_c \) & Number of continuous variables in \( T \). \\ \hline
\( \{C_1, \dots, C_{n_c}\} \) & Set of continuous random variables in \( T \). \\ \hline
\( n_d \) & Number of discrete variables in \( T \). \\ \hline
\( \{D_1, \dots, D_{n_d}\} \) & Set of discrete random variables in \( T \). \\ \hline
\( M(C_1, \dots, C_{n_c}, D_1, \dots, D_{n_d}) \) & Generative model used to learn the joint distribution $\rightarrow$ \( P(C_1, \dots, C_{n_c}, D_1, \dots, D_{n_d}) \). \\ \hline
\( P(C_1, \dots, C_{n_c}, D_1, \dots, D_{n_d}) \) & Joint distribution of continuous and discrete variables in \( T \). \\ \hline
\( T_{\text{pool}} \) & Pool of individuals generated by the model through statistical learning of the samples. \\ \hline
\( \mathcal{Z} \) & Set of zones. \\ \hline
\( z \) & Index of spatial zones in the target region. \\ \hline
\( h \) & Index of households in the pooled population. \\ \hline
\( \mathcal{I} \) & Set of individuals in the pooled dataset. \\ \hline
\( \mathcal{K} \) & Set of marginal tabulations. \\ \hline
\( k \) & Index of tabulations (characteristics or constraints) for which the population must be evaluated. \\ \hline
\( i \) & Index of cells within a tabulation. \\ \hline
\( x_{i,h} \) & Binary decision variable: 1 if individual \( i \) is assigned to household \( h \), 0 otherwise. \\ \hline
\( H_{h,z} \) &Count of people dwelling in house \( h \) in zone \( z \). \\ \hline
$M$ & Parameter that defines the max value a household can be filled \\ \hline
\( \mathcal{A} \) & Set of attributes (e.g., age, income, etc.). \\ \hline
\( \mathcal{P} \) & Set of all observed pairs of joint attributes (e.g., \( (D_1, D_2) \)) in the real dataset. \\ \hline
\( \mathcal{J} \) & Set of all joint attributes that define valid combinations of attributes. \\ \hline
\( J_{i,k} \) & Joint attribute \( k \) for individual \( i \). \\ \hline
\( \text{A}_{i,k} \) & Value of attribute \( k \) for individual \( i \). \\ \hline
\( P_z \) & Expected population at zone \( z \). \\ \hline
\( H_z \) & Expected households at zone \( z \). \\ \hline
\( O_{k,i} \) & Observed count for the \( i \)-th cell of the \( k \)-th tabulation. \\ \hline
\( E_{k,i} \) & Expected (known) count for the \( i \)-th cell of the \( k \)-th tabulation. \\ \hline
\( \text{RSSZ} \) & Relative Sum of Squared Z-scores, used to evaluate the fit of the synthetic population. \\ \hline
\( \text{RSSZ}_z \) & RSSZ value for spatial zone \( z \). \\ \hline
\( \text{SSZ}_{z,k} \) &  Sum of Squared Z-scores for tabulation \( k \) in zone \( z \). \\ \hline
\( N_k \) & Total count for tabulation \( k \). \\ \hline
\( C_k \) & Critical value of the chi-squared (\( \chi^2 \)) distribution for tabulation \( k \), adjusted for degrees of freedom. \\ \hline
\( \min \sum_{z} \text{RSSZ}_z \) & Objective function: minimize the cumulative RSSZ across all spatial zones. \\ \hline

\end{longtable}

\FloatBarrier
\subsection{Formulation}
This study addresses the challenge of synthesizing populations by ensuring that the distributions of the synthetic population are closely aligned with predefined constraining tabulations. The fit of the synthetic population to these tabulations is evaluated using the \gls{rssz} statistic, as proposed by \cite{Voas2001EvaluatingMicrodata} and further refined by \cite{Huang2001AMicrodata.}. The \gls{rssz} statistic provides a measure of goodness-of-fit between the synthetic population and the constraining tabulations because it takes into account the relative size of the category between actual and simulated populations. This ensures that categories with many levels are not given undue bias over smaller categories with fewer levels \citep{Rahman2023PopulationTechnique}. 
\\
The RSSZ statistic is calculated by adding the relative sum of squared Z scores (\( \text{RSSZ} \)) for each input tabulation \( k \), over the \(i\) categories of the k-th attribute. This formula for a zone, is expressed as follows:


\[
\text{RSSZ} = \frac{1}{N_k} \sum_{i} \left( \frac{O_{k,i} - E_{k,i}}{\sqrt{E_{k,i}}} \right)^2 \cdot \frac{1}{C_k}
\]

where:
\begin{itemize}
\item \( O_{k,i} \): Observed count for the \( i \)-th cell of the \( k \)-th tabulation.
\item \( E_{k,i} \): Expected (known) count for the \( i \)-th cell of the \( k \)-th tabulation.
\item \( N_k \): Total count for tabulation \( k \).
\item \( C_k \): Critical value of the chi-squared (\( \chi^2 \)) distribution for tabulation \( k \), adjusted for degrees of freedom (\( i-1 \), where \( i \) is the number of cells).

\end{itemize}
The use of \( C_k \) ensures that \gls{rssz} accounts for differences in the number of cells in the tabulations, allowing fair comparisons among tabulations of varying sizes.

The synthetic population is generated independently for each spatial zone using the simulated microdata. Since the simulated microdata includes SA4 information (which is a collection of SA2s), the SA2 generated independently can be utilized for zones that exist within the simulated microdata. The order in which the zones are processed does not impact the results. For each zone, a subset of individuals is iteratively refined to minimize the RSSZ) ensuring that the final synthetic population for each zone closely aligns with the constraining tabulations specific to that zone

\subsection*{Objective Function}
The objective function is to minimize the cumulative \gls{rssz} across all zones in the set \( \mathcal{Z} \). This is expressed as:

\[
\min \sum_{z \in \mathcal{Z}} \frac{\text{RSSZ}_z}{| \mathcal{Z} |}
\]
\\
where \( \text{RSSZ}_z \) is the relative Sum of Squared Z-scores (RSSZ) for zone \( z \).

This objective function minimizes deviations from the constraining tabulations, ensuring that the synthesized populations align closely with the true population distribution. The process is subject to the following constraints:

% \section*{Household Matching Constraints}

% For each household \( h \) in zone \( z \), an individual \( i \) can be part of household \( h \) only if the sum of the individuals' household attributes matches or is equal to the corresponding attributes of other individuals in the same household. This can be formulated as follows:

% Let:
% \begin{itemize}
%     \item \( \text{Attr}_{i,k} \): Attribute \( k \) of individual \( i \).
%     \item \( x_{i,h} \in \{0, 1\} \): Binary decision variable, where \( x_{i,h} = 1 \) if individual \( i \) is assigned to household \( h \), and \( 0 \) otherwise.
%     \item \( \text{AttrSum}_{h,k} \): Total value of attribute \( k \) for all individuals in household \( h \).
% \end{itemize}

\subsection*{Constraints}

\begin{enumerate}[label=\textbf{C\arabic*}, ref=C\arabic*]

    \item \textbf{Zone Population Constraints:}
    \begin{align}
        \sum_{h \in H_z} x_{i,h} \leq Z_z, \quad \forall z \in \mathcal{Z}
        \label{eq:zone_population}
    \end{align}  
    This ensures that the total population assigned to zone \( z \) does not exceed the total population \( Z_z \) of the zone.

    \item \textbf{Household Size Constraints:}
    \begin{align}
        \sum_{i \in \mathcal{I}} x_{i,h} = H_{h,z}, \quad \forall h \in \mathcal{H}_z, \, z \in \mathcal{Z}
        \label{eq:household_size}
    \end{align}
    Ensures that the number of individuals assigned to a household matches the predefined household size distribution.

  \item  \textbf{Marginal Attribute Matching:}
\begin{equation}
\sum_{h \in H} \sum_{\substack{i \in \mathcal{I} \\ (J_{i,k}, J_{i, k'}) \in \mathcal{P} \\ f(\mathcal{K}) =1 }} x_{i,h} \cdot A_{i,k} \leq \sum_{m \in \mathcal{K}} \text{tol}_{k,m} \cdot O_{k,m}, \quad \forall k \in \mathcal{K}
\label{eq:marginal_matching}
\end{equation}
    Matches the synthesized population assigned for all sample i assigned to a household to the marginal distributions of demographic attributes k over m levels, with respect to constraint rules, and sample observations.

    \item \textbf{Joint Distribution Consistency:}
\textbf{Joint Distribution Consistency:}
\begin{align}
\sum_{\substack{m \in \mathcal{K} }} \text{tol}_k \cdot O_{k,m} = E_{k,m}, \quad \forall k \in \mathcal{K}
        \label{eq:joint_distribution}
    \end{align}
    Ensures that the joint distributions align with observed counts \( O_{k,m} \) and expected counts \( E_{k,n} \) for each respective cell with respect to the tolerance level.

    \item \textbf{Joint Attribute Pair Constraint:}
\[
\mathcal{P} = \bigcup_{\substack{S \subseteq \{1, \dots, n_d\} \\ |S| = 2}} \{(D_k, D_{k'}) \mid \forall (k, k') \in S, \, (D_k^{(i)}, D_{k'}^{(i)}) \in T, \, \forall i \in \{1, \dots, N\}\}.
\]
This defines the valid pairs allowing:
    \[
    (J_{i,k} ,J_{i,k'}) \in \mathcal{P}, \quad \forall i \in \mathcal{I}, \, k, k' \in \mathcal{J}
    \label{eq:jatt}
    \]
    This ensures that all synthetic individuals follow valid joint attribute combinations observed in the real dataset.


    \item \textbf{Household Composition Constraints:}
    \begin{align}
        \sum_{i \in \mathcal{I}}\sum_{h \in \mathcal{H_z}} x_{i,h} \leq 1, \quad \forall z \in Z\\
         \sum_{i \in \mathcal{I}} x_{i,h} \leq M, \quad \forall h \in H_z z \in Z\\
        \label{eq:household_composition}
    \end{align}
   At most individual i can only be assigned to one household, and for each household a maximum number of individuals can be assigned

    \item \textbf{Non-Negativity and Binary Constraints:}
    \begin{align}
        tol_{i,h} > 0 & \forall i \in \mathcal{I}, \, h \in H
        \label{eq:binary_constraints}
    \end{align}
    Ensures that the assignment variables are binary (0 or 1).

    \item \textbf{Synthetic Population Consistency:}
    \begin{align}
        \sum_{z \in \mathcal{Z}} \sum_{h \in H_z} \sum_{i \in \mathcal{I}} x_{i,h} =  \sum_{z \in \mathcal{Z}} \sum_{k \in J} \sum_{i \in \mathcal{I}}E^z_{k,i}
        \label{eq:population_consistency}
    \end{align}
    Ensures that the total population across all zones matches the predefined total population.


\end{enumerate}
\section*{Filtering Function as as a Set of Rules}
In conjunction to the constraints \cref{eq:marginal_matching} uses a filtering function in the formulation.\\

The filtering function \( f(\{I\}) \) determines whether a combination of attributes satisfies a predefined set of rules \( \mathcal{R} \):

\[
f(I) =
\begin{cases}
1 & \text{if } \bigwedge_{j=1}^{m} R_j{(I)} \text{ is satisfied}, \\
0 & \text{otherwise.}
\end{cases}
\]

Where:
\begin{itemize}
    \item \( \mathcal{R} = \{R_1, R_2, \dots, R_m\} \) is the set of \( m \) rules.
    \item Each rule \( R_j \) specifies a condition that an Individual with defined characteristics \( I\) must satisfy.
\end{itemize}



Rules for the case study are applied in \cref{sec:casestudy} are outlined in \cref{sec:ref}.


\section*{Explanations of Constraints}

\begin{enumerate}[label=\textbf{C\arabic*}, ref=C\arabic*]

    \item (\cref{eq:zone_population}): Limits the total population assigned to a zone to its known population \( Z_z \).

    \item (\cref{eq:household_size}): Ensures that each household is assigned the correct number of individuals.

    \item (\cref{eq:marginal_matching}): Matches the synthesized population to the marginal distributions of demographic attributes (e.g., age, gender, income).

    \item (\cref{eq:joint_distribution}): Ensures that the joint distribution of attributes aligns with observed and expected counts.

    \item (\cref{eq:binary_constraints}): Ensures that the observed and predicted have a tolerance level that is positive.

    \item (\cref{eq:population_consistency}): Ensures that the total population matches the predefined total.





% \section*{Additional Constraints}
% In addition to the above constraints, there are several logical constraints intrinsically derived from the relationships within the data, or provided by the rules as defined by the user. For further details, see \cref{sec:ref}.

\end{enumerate}
